\documentclass[12pt,a4paper]{article}

% Paquetes básicos
\usepackage[utf8]{inputenc}
\usepackage[T1]{fontenc}
\usepackage[spanish]{babel}
\usepackage{geometry}
\geometry{margin=2.5cm}
\usepackage{amsmath,amssymb,amsfonts}
\usepackage{graphicx}
\usepackage{booktabs}
\usepackage{hyperref}
\usepackage{fancyhdr}
\usepackage{titlesec}
\usepackage{enumitem}
\usepackage{xcolor}

% Configuración de hipervínculos
\hypersetup{
    colorlinks=true,
    linkcolor=blue!70!black,
    urlcolor=blue!80!black,
    citecolor=green!60!black
}

% Encabezado y pie
\pagestyle{fancy}
\fancyhf{}
\fancyhead[L]{Revisión Física e Ingeniería -- Simulador Multirotor}
\fancyhead[R]{\thepage}
\renewcommand{\headrulewidth}{0.4pt}

% Título de secciones
\titleformat{\section}{\Large\bfseries\color{blue!70!black}}{\thesection}{1em}{}
\titleformat{\subsection}{\large\bfseries\color{blue!50!black}}{\thesubsection}{1em}{}

\title{\textbf{Revisión de Extremo a Extremo:\\Física, Ingeniería y Matemáticas del Simulador Multirotor}}
\author{Análisis extraído del repositorio\\ \texttt{github.com/ChristianGonper/modelado-control-dron}}
\date{\today}

\begin{document}

\maketitle

\begin{abstract}
Este informe presenta una revisión exhaustiva y estructurada de los modelos físicos, matemáticos e ingenieriles implementados en el simulador multirotor del repositorio \texttt{ChristianGonper/modelado-control-dron}. 

Se extraen y documentan de forma clara las ecuaciones de dinámica de cuerpo rígido 6DOF, cinemática con cuaterniones, modelo de rotores y actuadores de primer orden, mezclador por pseudo-inversa, perturbaciones aerodinámicas (arrastre, viento Ornstein-Uhlenbeck e inducido), integración numérica RK4 y las bases temporales desacopladas. 

El objetivo es proporcionar una documentación autocontenida y reutilizable para rehacer o extender el modelado físico, separando completamente la lógica de ingeniería de cualquier concepto de arquitectura de software.
\end{abstract}

\newpage
\tableofcontents
\newpage

\section{Introducción y Alcance}

El simulador modela un vehículo multirotor como un cuerpo rígido con seis grados de libertad acoplados. La revisión se centra exclusivamente en:

\begin{itemize}
    \item Marcos de referencia y representación del estado.
    \item Ecuaciones de cinemática y dinámica.
    \item Modelo de actuadores y mezclador.
    \item Perturbaciones aerodinámicas y viento.
    \item Métodos numéricos de integración y discretización temporal.
    \item Límites físicos y suposiciones de modelado.
\end{itemize}

Todo el contenido ha sido extraído directamente de los documentos de teoría (\texttt{docs/theory/}), decisiones arquitectónicas relevantes (ADRs) y código fuente de los módulos \texttt{core/} y \texttt{dynamics/} mediante la GitHub REST API.

\section{Marcos de Referencia y Representación del Estado}

\subsection{Estado del Vehículo (\texttt{VehicleState})}

El estado se define mediante:

\begin{itemize}
    \item Posición $\mathbf{p} \in \mathbb{R}^3$ y velocidad lineal $\mathbf{v} \in \mathbb{R}^3$ en \textbf{marco inercial (mundo)}.
    \item Orientación representada por cuaternión unitario $\mathbf{q} = (w, x, y, z)$.
    \item Velocidad angular $\boldsymbol{\omega} \in \mathbb{R}^3$ en \textbf{marco cuerpo}.
\end{itemize}

\textbf{Validaciones implementadas}:
\begin{itemize}
    \item $\mathbf{q}$ debe ser unitario: $\|\mathbf{q}\| = 1$ (con tolerancia $10^{-6}$).
    \item Todas las componentes deben ser finitas y reales.
\end{itemize}

\subsection{Cinemática con Cuaterniones}

Las funciones principales implementadas en \texttt{core/attitude.py} son:

\begin{align}
    \dot{\mathbf{q}} &= \frac{1}{2} \mathbf{q} \otimes [0, \boldsymbol{\omega}] \\
    \mathbf{q}_{\text{err}} &= \text{normalize}(\mathbf{q}_{\text{des}} \otimes \mathbf{q}_{\text{conj}}) \\
    \boldsymbol{\theta}_{\text{vec}} &= 2 \cdot \text{atan2}(\|\mathbf{v}\|, w) \cdot \frac{\mathbf{v}}{\|\mathbf{v}\|}
\end{align}

donde $\otimes$ denota multiplicación de cuaterniones y $\mathbf{v} = (x,y,z)$ del error.

\textbf{Ventaja}: evita singularidades de los ángulos de Euler y permite una propagación geométrica estable.

\section{Dinámica de Cuerpo Rígido 6DOF}

\subsection{Ecuaciones de Movimiento}

\textbf{Dinámica translacional} (marco mundo):
\begin{equation}
    m \frac{d\mathbf{v}}{dt} = \mathbf{R}(\mathbf{q}) \cdot \mathbf{T}_{\text{body}} + m\mathbf{g} + \mathbf{F}_{\text{aero}} + \mathbf{F}_{\text{viento}}
\end{equation}

\textbf{Dinámica rotacional} (marco cuerpo):
\begin{equation}
    \mathbf{I} \frac{d\boldsymbol{\omega}}{dt} = \boldsymbol{\tau}_{\text{body}} - \boldsymbol{\omega} \times (\mathbf{I} \boldsymbol{\omega})
\end{equation}

donde $\mathbf{I}$ es el tensor de inercia (asumido diagonal).

\subsection{Suposiciones del Modelo}

\begin{itemize}
    \item Cuerpo rígido único con masa concentrada e inercia diagonal.
    \item Sin flexibilidad estructural ni flapping de rotores.
    \item Sin modelo de batería ni efectos de contacto.
    \item Empuje y par agregados a nivel vehículo (aunque internamente se resuelven por rotor).
\end{itemize}

\section{Modelo de Rotores y Actuación}

\subsection{Parámetros por Rotor (\texttt{RotorGeometry})}

\begin{itemize}
    \item Posición en marco cuerpo $\mathbf{r}_i \in \mathbb{R}^3$.
    \item Dirección del eje (normalmente $(0,0,1)$).
    \item Sentido de giro $\text{spin\_direction} = \pm 1$.
    \item Coeficiente de empuje $k_T$ [N/(rad/s)$^2$].
    \item Coeficiente de par de reacción $k_Q$ [Nm/N].
    \item Constante de tiempo del motor $\tau_m$ [s].
    \item Velocidad angular máxima $\omega_{\max}$ [rad/s].
\end{itemize}

\subsection{Dinámica del Actuador (Primer Orden)}

\begin{equation}
    \omega_{\text{target}} = \sqrt{\frac{T_{\text{target}}}{k_T}}
\end{equation}

\begin{equation}
    \alpha = 1 - e^{-\Delta t / \tau_m}, \quad
    \omega_{t+1} = \omega_t + \alpha (\omega_{\text{target}} - \omega_t)
\end{equation}

\begin{equation}
    T_{t+1} = k_T \cdot \omega_{t+1}^2
\end{equation}

\textbf{Saturación}: $\omega_{t+1} \leftarrow \text{clamp}(\omega_{t+1}, 0, \omega_{\max})$.

\textbf{Torque de reacción}:
\begin{equation}
    \tau_{\text{reac},i} = T_i \cdot k_Q \cdot \text{spin\_direction}_i
\end{equation}

\section{Mezclador y Asignación de Empuje}

El controlador genera un \texttt{VehicleIntent} = $(T_{\text{colectivo}}, \boldsymbol{\tau}_{\text{body}})$.

El mezclador resuelve:
\begin{equation}
    \mathbf{T} = \arg\min_{\mathbf{T} \geq 0} \left\| \mathbf{A} \mathbf{T} - \begin{bmatrix} T_{\text{coll}} \\ \tau_x \\ \tau_y \\ \tau_z \end{bmatrix} \right\|^2_2
\end{equation}

donde las columnas de $\mathbf{A}$ incorporan la posición del rotor y el coeficiente $k_Q$.

Posteriormente se aplica saturación individual por rotor:
\begin{equation}
    T_i \leftarrow \min(T_i, k_T \cdot \omega_{\max}^2)
\end{equation}

\section{Aerodinámica y Perturbaciones}

\subsection{Arrastre Parásito}

\begin{equation}
    \mathbf{F}_{\text{drag}} = -\frac{1}{2} \rho A_{\text{drag}} \|\mathbf{v}_{\text{rel}}\| \mathbf{v}_{\text{rel}}
\end{equation}

(activable mediante parámetro; $\rho \approx 1.225$ kg/m$^3$ por defecto).

\subsection{Pérdida de Eficiencia en Hover Inducida}

Reducción proporcional del empuje efectivo cuando el comando se acerca al máximo (parámetro \texttt{induced\_hover\_loss\_ratio}).

\subsection{Modelo de Viento (Ornstein-Uhlenbeck)}

Componente base determinista + proceso estocástico dt-invariante:

\begin{equation}
    g_{t+1} = \mu + e^{-\Delta t / \tau_g} (g_t - \mu) + \sigma \sqrt{1 - e^{-2\Delta t / \tau_g}} \cdot \mathcal{N}(0,1)
\end{equation}

\textbf{Ventaja}: la varianza estacionaria es independiente del paso de tiempo $\Delta t$.

\subsection{Observabilidad}

Se distingue explícitamente entre:
\begin{itemize}
    \item \texttt{true\_state}: estado físico real de la planta.
    \item \texttt{observed\_state}: estado expuesto al controlador (con perturbaciones).
    \item \texttt{tracking\_state}: estado usado para métricas de error (por defecto = true).
\end{itemize}

\section{Integración Numérica y Bases Temporales}

\subsection{Método de Integración}

La planta se integra con \textbf{RK4} (Runge-Kutta de orden 4) en pasos de \texttt{physics\_dt\_s}.

\subsection{Desacoplamiento de Frecuencias}

\begin{itemize}
    \item \texttt{physics\_dt\_s}: avance de la dinámica rígida.
    \item \texttt{control\_dt\_s}: frecuencia de cálculo del \texttt{VehicleIntent}.
    \item \texttt{telemetry\_dt\_s}: muestreo y exportación de historia.
\end{itemize}

Este desacoplamiento permite estudiar efectos de lag de actuador y muestreo sin modificar el contrato del controlador.

\section{Parámetros Físicos y Límites}

El simulador está motivado por la plataforma \textbf{LiteWing} (cuadricóptero ligero con motores coreless DC), pero \textbf{no utiliza parámetros identificados de vuelo real}.

Parámetros configurables principales:
\begin{itemize}
    \item Densidad del aire: $\rho = 1.225$ kg/m$^3$.
    \item Área de arrastre $A_{\text{drag}}$.
    \item Coeficientes de rotores ($k_T$, $k_Q$, $\tau_m$, $\omega_{\max}$).
    \item Masa e inercia (definidos por escenario).
    \item Intensidad y tiempo de correlación del viento.
\end{itemize}

\textbf{Límites de validez}:
\begin{itemize}
    \item Modelo válido para experimentos de control y robustez controlados.
    \item No constituye un gemelo digital de alta fidelidad.
    \item No incluye flapping, dinámica de batería ni estimador onboard.
\end{itemize}

\section{Resumen de Fórmulas Clave}

\begin{table}[h]
\centering
\begin{tabular}{@{}ll@{}}
\toprule
\textbf{Dominio} & \textbf{Fórmula Principal} \\
\midrule
Cinemática & $\dot{\mathbf{q}} = \frac12 \mathbf{q} \otimes [0, \boldsymbol{\omega}]$ \\
Dinámica translacional & $m \dot{\mathbf{v}} = \mathbf{R}(\mathbf{q})\mathbf{T} + m\mathbf{g} + \mathbf{F}_{\text{aero}}$ \\
Dinámica rotacional & $\mathbf{I} \dot{\boldsymbol{\omega}} = \boldsymbol{\tau} - \boldsymbol{\omega} \times (\mathbf{I}\boldsymbol{\omega})$ \\
Actuador motor & $\omega_{t+1} = \omega_t + (1-e^{-\Delta t/\tau})(\omega_{\text{target}}-\omega_t)$ \\
Mezclador & $\mathbf{T} = \arg\min \|\mathbf{A}\mathbf{T} - \mathbf{b}\|_2^2$ (con saturación) \\
Arrastre & $\mathbf{F}_d = -\frac12\rho A \|\mathbf{v}\| \mathbf{v}$ \\
Viento OU & $g_{t+1} = \mu + e^{-\Delta t/\tau}(g_t-\mu) + \sigma\sqrt{1-e^{-2\Delta t/\tau}}\mathcal{N}(0,1)$ \\
Error de actitud & $\boldsymbol{\theta}_{\text{vec}} = 2 \cdot \text{atan2}(\|\mathbf{v}\|,w) \cdot \frac{\mathbf{v}}{\|\mathbf{v}\|}$ \\
\bottomrule
\end{tabular}
\caption{Resumen de ecuaciones principales del simulador}
\end{table}

\section{Limitaciones y Recomendaciones para Rehacer el Modelo}

\textbf{Fortalezas actuales}:
\begin{itemize}
    \item Coherencia física para experimentos de TFG.
    \item Cuaterniones + RK4 + actuador de primer orden + viento OU = excelente equilibrio realismo/computacional.
    \item Separación clara entre \textit{intent} del controlador e \textit{acción aplicada}.
\end{itemize}

\textbf{Limitaciones importantes}:
\begin{itemize}
    \item Sin modelo de flapping ni aerodinámica por rotor.
    \item Sin dinámica de batería ni sag de voltaje.
    \item Sin estimador onboard, delays de sensores ni fusión.
    \item Parámetros no identificados de vuelo real (solo ``LiteWing-inspired'').
    \item Integración síncrona (no event-driven ni multi-rate de firmware real).
\end{itemize}

\textbf{Recomendaciones para una versión mejorada}:
\begin{enumerate}
    \item Incorporar modelo de flapping o inflow de rotor.
    \item Añadir dinámica simple de batería.
    \item Hacer el mezclador más general (soporte hexacóptero, rotores inclinados, etc.).
    \item Validar contra telemetría real del LiteWing.
    \item Publicar el conjunto de parámetros identificados.
\end{enumerate}

\section{Control: Estrategias PD en Cascada y Redes Neuronales}

\subsection{Contrato Común de Controladores}

Todos los controladores (PD clásico y neuronales) implementan la interfaz \texttt{ControllerContract}:

\begin{itemize}
    \item \texttt{kind}: tipo de controlador (``cascade'', ``mlp'', ``gru'', ``lstm'', etc.).
    \item \texttt{compute\_action(observation, reference)} $\to$ \texttt{VehicleCommand}.
\end{itemize}

Esto permite intercambiar controladores sin modificar el runner ni la planta.

\subsection{Control PD en Cascada (Fase 1 -- Control Clásico)}

El controlador en cascada (\texttt{cascade.py}) sigue la arquitectura clásica de dos bucles:

\subsubsection{Bucle Externo (Posición)}
\begin{align}
    \mathbf{e}_p &= \mathbf{p}_{\text{ref}} - \mathbf{p} \\
    \mathbf{e}_v &= \mathbf{v}_{\text{ref}} - \mathbf{v} \\
    \mathbf{a}_{\text{des}} &= \mathbf{K}_p \mathbf{e}_p + \mathbf{K}_d \mathbf{e}_v + \mathbf{a}_{\text{ff}}
\end{align}

Donde $\mathbf{K}_p = (2.0, 2.0, 4.0)$ y $\mathbf{K}_d = (1.2, 1.2, 2.5)$ por defecto (ganancias sintonizadas para el escenario nominal).

A partir de $\mathbf{a}_{\text{des}}$ se calcula el empuje colectivo deseado y la actitud deseada:
\begin{equation}
    \mathbf{F}_{\text{des}} = m (\mathbf{a}_{\text{des}} + g \mathbf{e}_z), \quad
    T_{\text{coll}} = \|\mathbf{F}_{\text{des}}\|
\end{equation}

La actitud deseada se obtiene con \texttt{quaternion\_from\_thrust\_direction\_and\_yaw}.

\subsubsection{Bucle Interno (Actitud)}
\begin{align}
    \boldsymbol{\theta}_{\text{err}} &= \text{quaternion\_error\_vector}(\mathbf{q}_{\text{des}}, \mathbf{q}) \\
    \boldsymbol{\tau}_{\text{cmd}} &= \mathbf{K}_{p,\text{att}} \boldsymbol{\theta}_{\text{err}} + \mathbf{K}_{d,\text{rate}} (\boldsymbol{\omega}_{\text{des}} - \boldsymbol{\omega})
\end{align}

con saturación por eje según \texttt{max\_body\_torque\_nm} y \texttt{max\_tilt\_rad} = 0.6 rad (≈34°).

\textbf{Parámetros por defecto} (sintonizados para el escenario mínimo):
- Bucle posición: $K_p = [2.0, 2.0, 4.0]$, $K_d = [1.2, 1.2, 2.5]$
- Bucle actitud: $K_p = [4.5, 4.5, 1.8]$, $K_d = [0.15, 0.15, 0.08]$

\subsection{Controladores Neuronales (Fases 3 y 4)}

\subsubsection{Arquitectura MLP (Fase 3 -- Feedforward)}

\begin{itemize}
    \item Entrada: ventana de características (\texttt{window\_size} × \texttt{feature\_dim}).
    \item Características: ``raw\_observation'' o ``observation\_plus\_tracking\_errors''.
    \item Capas ocultas: (128, 64) con ReLU.
    \item Salida: 4 acciones (empuje colectivo + 3 componentes de torque).
\end{itemize}

\textbf{Entrenamiento}:
\begin{itemize}
    \item Ventanas deslizantes con \texttt{stride}.
    \item Normalización (media y desviación estándar calculadas solo en train).
    \item Pérdida MSE sobre acciones.
    \item Checkpoint guarda: arquitectura + estadísticas de normalización + pesos.
\end{itemize}

\subsubsection{Arquitectura Recurrente (Fase 4 -- GRU / LSTM)}

\begin{itemize}
    \item Mantiene estado oculto entre pasos (memoria temporal).
    \item Arquitecturas soportadas: GRU y LSTM.
    \item Tamaño de ventana por defecto: 8 (GRU) / 16 (LSTM).
    \item Mismo pipeline de características y normalización que el MLP.
\end{itemize}

\textbf{Ventaja teórica}: captura dinámicas temporales y dependencias de largo plazo que un MLP estático no puede modelar.

\subsubsection{Teoría Seguida para los Controladores Neuronales}

\begin{enumerate}
    \item \textbf{Feature Engineering}: se usan errores de seguimiento (posición, velocidad, actitud) como características adicionales para que la red aprenda directamente la corrección de error.
    \item \textbf{Normalización por dataset}: media y std calculadas exclusivamente en el conjunto de entrenamiento para evitar fuga de datos.
    \item \textbf{Reproducibilidad}: semillas fijas para Python, NumPy y PyTorch + algoritmos determinísticos cuando es posible.
    \item \textbf{Intercambiabilidad}: todos los modelos neuronales devuelven exactamente el mismo tipo \texttt{VehicleCommand} que el controlador PD, permitiendo comparación justa en los mismos escenarios.
    \item \textbf{Checkpointing robusto}: se persiste no solo los pesos, sino también la arquitectura y las estadísticas de normalización para garantizar inferencia idéntica en producción.
\end{enumerate}

\section{Conclusiones}

Este informe proporciona una documentación clara, autocontenida y lista para reutilizar de toda la lógica física e ingenieril del simulador multirotor, incluyendo ahora las estrategias de control clásico (PD en cascada) y las aproximaciones neuronales (MLP y recurrentes).

La separación explícita entre la física (este documento) y la arquitectura de software (futura revisión) facilita tanto la extensión del modelo como la implementación de nuevos controladores (clásicos o neuronales) sobre una base física sólida.

\vspace{1cm}
\begin{center}
\textit{Documento generado automáticamente a partir del análisis del repositorio}\\
\texttt{https://github.com/ChristianGonper/modelado-control-dron}\\
\textit{Fecha: \today}
\end{center}

\end{document}